\documentclass[11pt]{article}
\usepackage[margin=1in]{geometry}
\usepackage{amsmath, amssymb}

\title{Exercise 3.14: Verifying the Bellman Equation for the Gridworld}
\author{Sutton \& Barto, Section 3.5}
\date{}

\begin{document}
\maketitle

\section*{Problem}

The Bellman equation (3.14) must hold for each state for the value function
$v_\pi$ shown in Figure~3.2 (right) of Example~3.5. Show numerically that
this equation holds for the center state, valued at $+0.7$, with respect to
its four neighboring states, valued at $+2.3$, $+0.4$, $-0.4$, and $+0.7$.

\section*{Solution}

\subsection*{Setup}

The Bellman equation (3.14) is:
\[
    v_\pi(s) = \sum_{a} \pi(a \mid s) \sum_{s'} \sum_{r}
    p(s', r \mid s, a) \left[ r + \gamma\, v_\pi(s') \right]
\]

For the equiprobable random policy, $\pi(a \mid s) = 0.25$ for each of the
four actions. The center state is not a special state (not $A$ or $B$), so each
action deterministically moves to the neighboring cell with reward $r = 0$.
The discount factor is $\gamma = 0.9$.

\subsection*{Neighboring values}

\begin{itemize}
    \item North: $v_\pi = +2.3$
    \item South: $v_\pi = -0.4$
    \item East: $v_\pi = +0.4$
    \item West: $v_\pi = +0.7$
\end{itemize}

\subsection*{Verification}

Since the transitions are deterministic and all rewards are 0:
\begin{align*}
    v_\pi(s) &= \sum_{a} \pi(a \mid s) \left[ r + \gamma\, v_\pi(s') \right] \\[6pt]
             &= 0.25 \left[ 0 + 0.9 \cdot 2.3 \right]
              + 0.25 \left[ 0 + 0.9 \cdot (-0.4) \right] \\
             &\quad + 0.25 \left[ 0 + 0.9 \cdot 0.4 \right]
              + 0.25 \left[ 0 + 0.9 \cdot 0.7 \right] \\[6pt]
             &= 0.25 \left[ 2.07 + (-0.36) + 0.36 + 0.63 \right] \\[4pt]
             &= 0.25 \times 2.7 \\[4pt]
             &= 0.675 \\[4pt]
             &\approx \boxed{0.7} \;\checkmark
\end{align*}

This confirms the Bellman equation holds (to one decimal place) for the center state.

\end{document}
